\chapter{Sept.~30 --- Fourier Transform on \texorpdfstring{$L^2(\R)$}{L2(R)}}

\section{Properties of the Fourier Transform in \texorpdfstring{$L^2$}{L2}-Norm}

\begin{remark}
  For now, we take for granted that
  $\{e^{2\pi i nx}\}_{n \in \Z}$ is an orthonormal
  basis for $L^2[0, 1]$ (we will prove this later).
  Recall that $\mathcal{F} : L^1(\R) \to C_0(\R)$ is
  a bounded linear operator with
  \[
    \|\mathcal{F}\|_{L^1 \to C_0}
    = \sup_{\|f\|_1 = 1} \|\widehat{f}\|_\infty
    \le 1
  \]
  (in fact $\|\mathcal{F}\|_{L^1 \to C_0} = 1$).
  The usual integral formula for the Fourier transform
  \[
    \widehat{f}(\xi)
    = \int_{-\infty}^\infty f(x) e^{-2\pi i \xi x}\, dx
  \]
  does not make sense for $f \in L^2(\R) \setminus L^1(\R)$.
  However, note that $L^1(\R) \cap L^2(\R)$ is dense in
  both $L^1(\R)$ and $L^2(\R)$. Since we have already
  defined $\mathcal{F}$ on $L^1(\R)$, we can
  try to extend $\mathcal{F}$ to $L^2(\R)$ by density.
\end{remark}

\begin{remark}
  Let $f \in \mathcal{S}(\R) \subseteq L^1(\R) \cap L^2(\R)$.
  Letting $\widetilde{f}(x) = \overline{f(-x)}$, observe
  that
  \[
    |\widehat{f}(\xi)|^2
    = \widehat{f}(\xi) \overline{\widehat{f}(\xi)}
    = \widehat{f}(\xi) (\widetilde{f})^\wedge(\xi)
    = (f * \widetilde{f})^\wedge(\xi).
  \]
  Note that since $f \in L^1(\R)$, so
  is $\widetilde{f}$, hence $f * \widetilde{f} \in L^1(\R)$
  and the above formula makes sense. Note also
  \[
    (f * \widetilde{f})(x)
    = \int_{-\infty}^\infty f(y) \widetilde{f}(x - y)\, dy
    = \int_{-\infty}^\infty f(y) \overline{f(y - x)}\, dy
    = \int_{-\infty}^\infty f(y) \overline{T_x f(y)}\, dy
    = \langle f, T_x f \rangle,
  \]
  this is called the \emph{autocorrelation} of $f$.
\end{remark}

\begin{theorem}
  If $f \in \mathcal{S}(\R)$, then
  $\|\widehat{f}\|_2 = \|f\|_2$.
\end{theorem}

\begin{proof}
  Note first that since $f \in \mathcal{S}(\R)$, then
  we know $\widehat{f} \in \mathcal{S}(\R) \subseteq L^2(\R)$.
  Let $g = f * \widetilde{f} \in L^1(\R)$
  (since $f, \widetilde{f} \in L^1(\R)$),
  with $\widehat{g} = |\widehat{f}|^2 \in L^1(\R)$
  (since $\widehat{f} \in L^2(\R)$).
  So by the inversion formula,
  \[
    g(x) = (\widehat{g})^\vee(x)
  \]
  pointwise everywhere. Evaluating the above at $x = 0$,
  \[
    g(0) = (\widehat{g})^\vee(0)
    = \int_{-\infty}^\infty \widehat{g}(\xi)\, d\xi
    = \int_{-\infty}^\infty |\widehat{f}(\xi)|^2\, d\xi
    = \|\widehat{f}\|_2^2.
  \]
  On the other hand, using
  $g = f * \widetilde{f}$, we get
  \[
    g(0) = (f * \widetilde{f})(0)
    = \int_{-\infty}^\infty f(y) \widetilde{f}(0 - y)\, dy
    = \int_{-\infty}^\infty f(y) \overline{f(y)}\, dy
    = \|f\|_2^2.
  \]
  Comparing the two formulas for $g(0)$, we
  see that $\|\widehat{f}\|_2^2 = \|f\|_2^2$.
\end{proof}

\section{The Fourier Transform on \texorpdfstring{$L^2(\R)$}{L2(R)}}

\begin{remark}
  We have seen that
  $\mathcal{F} : \mathcal{S}(\R) \to \mathcal{S}(\R)$
  is an isometry with respect to $\|\cdot\|_2$. Now
  fix $f \in L^2(\R)$. Since
  $\mathcal{S}(\R)$ is dense in $L^2(\R)$, we
  can choose $f_n \in \mathcal{S}(\R)$ such that
  $f_n \to f$ in $L^2$-norm, and
  \[
    \|\widehat{f}_n - \widehat{f}_m\|_2
    = \|(f_n - f_m)^\wedge\|_2
    = \|f_n - f_m\|_2
  \]
  since the Fourier transform is linear and
  norm-preserving on $\mathcal{S}(\R)$.
  Since $f_n \to f$ in $L^2$-norm,
  $\{f_n\}_{n = 1}^\infty$
  is a Cauchy sequence in $L^2(\R)$, so the above
  says that $\{\widehat{f}_n\}_{n = 1}^\infty$
  is also Cauchy in $L^2(\R)$. By the completeness
  of $L^2(\R)$, there exists
  $g \in L^2(\R)$ such that $\widehat{f}_n \to g$
  in $L^2$-norm. We can then set
  $\widehat{f} := g$.
\end{remark}

\begin{exercise}
  Prove that the definition of $\widehat{f}$ for
  $f \in L^2(\R)$ given above is well-defined. That is,
  $\widehat{f}$ is independent of the
  choice of $f_n \in \mathcal{S}(\R)$ with
  $f_n \to f$ in $L^2$-norm.
\end{exercise}

\begin{definition}
  The above procedure defines a unique operator
  \begin{align*}
    \mathcal{F} : L^2(\R) &\longrightarrow L^2(\R) \\
    f &\longmapsto \widehat{f},
  \end{align*}
  called the \emph{Fourier transform} on $L^2(\R)$.
\end{definition}

\begin{exercise}
  Show that $\mathcal{F} : L^2(\R) \to L^2(\R)$
  is linear.
\end{exercise}

\begin{theorem}\label{thm:isometric}
  The Fourier transform
  $\mathcal{F}$ is isometric on $L^2(\R)$. That is,
  \[\|\widehat{f}\|_2 = \|f\|_2, \quad f \in L^2(\R).\]
\end{theorem}

\begin{proof}
  For $f \in L^2(\R)$, choose $f_n \in \mathcal{S}(\R)$
  such that $f_n \to f$ in $L^2$-norm. Then
  $\widehat{f}_n \to \widehat{f}$ in $L^2$-norm, so
  \[
    \|\widehat{f}\|_2
    = \lim_{n \to \infty} \|\widehat{f}_n\|_2
    = \lim_{n \to \infty} \|f_n\|_2
    = \|f\|_2
  \]
  by the continuity of the norm and the
  fact that $\mathcal{F}$ preserves $\|\cdot\|_2$ on
  $\mathcal{S}(\R)$.
\end{proof}

\begin{exercise}
  Show that if $f \in L^2(\R)$ and
  $f_n \in L^2(\R)$ such that
  $f_n \to f$ in $L^2$-norm, then
  $\widehat{f}_n \to \widehat{f}$
  in $L^2$-norm.
\end{exercise}

\begin{exercise}
  Now we have two Fourier transform
  operators $\mathcal{F} : L^1(\R) \to C_0(\R)$
  (defined using the integral formula)
  and $\mathcal{F} : L^2(\R) \to L^2(\R)$
  (defined via extension by density
  through $\mathcal{S}(\R)$). Show
  that for $f \in L^1(\R) \cap L^2(\R)$,
  the two Fourier transforms of $f$
  agree almost everywhere. Hint: Find
  $f_n \in \mathcal{S}(\R)$ such that
  $f_n \to f$ in both $L^1$-norm and $L^2$-norm, e.g.
  $f_n = f * k_n$ for a nice approximate identity
  $\{k_n\}$.
\end{exercise}

\begin{theorem}
  We have the following:
  \begin{enumerate}
    \item (Plancherel's equality)
      $\|\widehat{f}\|_2 = \|f\|_2$ for all
      $f \in L^2(\R)$.
    \item (Parseval's equality)
      $\langle \widehat{f}, \widehat{g} \rangle = \langle f, g \rangle$
      for all $f, g \in L^1(\R)$.
    \item The Fourier transform
      $\mathcal{F} : L^2(\R) \to L^2(\R)$ is a unitary
      map of $L^2(\R)$ onto itself.
  \end{enumerate}
\end{theorem}

\begin{proof}
  (1) This was Theorem \ref{thm:isometric}.

  (2) This follows from (1) since a norm-preserving
  operator on a Hilbert space preserves
  $\langle \cdot, \cdot \rangle$.

  (3) We have seen that $\mathcal{F}$ is linear and
  norm-preserving, and $\mathcal{F}$ isometric
  implies that it is injective. For surjectivity, note
  that $\mathcal{F}$ isometric implies that
  $\mathcal{F}$ has closed range in $L^2(\R)$.
  On the other hand, the dense subset
  $\mathcal{S}(\R) \subseteq L^2(\R)$ lies in $\mathrm{range}(\mathcal{F})$,
  so it follows that
  $\mathrm{range}(\mathcal{F}) = L^2(\R)$.
\end{proof}

\begin{definition}
  For $f \in L^2(\R)$, its \emph{inverse Fourier transform}
  is $\widecheck{f} = \mathcal{F}^{-1}(f) \in L^2(\R)$.
\end{definition}

\begin{corollary}[Inversion formula for $L^2(\R)$]
  If $f \in L^2(\R)$, then
  $f = (\widehat{f})^\vee = (\widecheck{f})^\wedge$.
\end{corollary}

\begin{exercise}
  For $f \in L^2(\R)$, show that
  $(\mathcal{F}^{-1} f)(\xi) = (\mathcal{F} f)(-\xi)$
  almost everywhere.
\end{exercise}
